% Recuperación expositiva desde los propietarios del núcleo TPK.
% La procedencia y el control de conservación se documentan en technical/TPK_INTEGRACION.md.
% Este archivo amplía la sección TPK; no sustituye extension.tex ni generacion.tex.

\subsubsection{Estado dinámico, observables y dependencia de las operaciones}
\label{tpk-int:estado}

El núcleo topológico de fase organiza la evolución de las dos hojas de APP
con la orientación determinada por TRIT. Su definición comprende el estado
sobre el que actúa, las reglas de transición y las relaciones que hacen
compatibles las prolongaciones. La ecuación \eqref{eq:actualizacion} adquiere
contenido operativo al especificar qué coordenadas recibe y modifica cada
uno de sus factores. La posición pertenece al soporte toroidal; la fase
pertenece al calendario; la orientación determina el transporte; los
cocientes registran la evaluación aritmética anterior a su reducción.
Estas coordenadas se relacionan mediante la transición y conservan funciones
matemáticas diferentes.

La proyección observable de un estado es
\begin{equation}
 q=(x_+,d_+,x_\times,d_\times,\theta)
 \in\mathcal Q_{\rm obs}
 :=(X_9\times\mathcal D)^2\times\Theta_{108},
 \qquad \mathcal D=\{N,E,S,O\},\quad
 \Theta_{108}=\ZZ/108\ZZ.
 \label{tpk-int:qobs}
\end{equation}
Los subíndices distinguen los cursores aditivo y multiplicativo. En las
fibras de estados completos sobre \(q\) se conservan la trayectoria ordenada, las evaluaciones
enteras de ambas hojas, sus residuos y cocientes, y los registros de
orientación y frontera. En una prolongación de precisión se añaden el
prefijo y el cilindro compatible. Por tanto, dos estados con la misma
proyección \(q\) pueden conservar registros distintos.

La dependencia de las operaciones tiene tres niveles. En el nivel local,
el selector activa una hoja, el transporte modifica su posición y la
actualización calcula las nuevas coordenadas aritméticas. En el nivel de
trayectorias, estas transiciones se componen y conservan los registros de
los pasos anteriores. En el nivel de refinamiento, las relaciones de
extensión determinan qué continuación conserva el prefijo, la firma y la
frontera del antecedente. El calendario finito describe la fase de esas
operaciones; la holonomía describe su composición sobre la historia.

\subsubsection{Selección trítica y transporte cardinal levantado}
\label{tpk-int:seleccion}

Sea \(\bar\theta\in\{0,\ldots,107\}\) el representante del calendario.
La fase operativa y el sector dodecafásico son
\begin{equation}
 r(\theta)=1+(\bar\theta\bmod9),\qquad
 \beta(\theta)=\left[\left\lfloor\frac{\bar\theta}{9}\right\rfloor\right]_{12}.
 \label{tpk-int:fase}
\end{equation}
El selector ya definido puede escribirse \(\varepsilon_9=M_{\rm ph}\).
Su vector de valores, en la carta ordenada de \(D_9\), es
\[
 m_{\rm ph}=(1,-1,0,1,-1,0,1,-1,0).
\]
La función \(M_{\rm ph}:D_9\to\TRIT\) y este vector son dos
presentaciones del mismo selector, con sus dominios respectivos.

Definamos los indicadores de activación
\[
 a_+(\theta)=\mathbf1_{\{1,4,7\}}(r(\theta)),\qquad
 a_\times(\theta)=\mathbf1_{\{2,5,8\}}(r(\theta)).
\]
El transporte de los cursores es entonces
\begin{equation}
 x_+'=x_++a_+(\theta)\nu_{d_+},\qquad
 x_\times'=x_\times+a_\times(\theta)\nu_{d_\times}
 \quad\text{en }X_9.
 \label{tpk-int:cursores}
\end{equation}
En las fases \(3,6,9\) permanecen ambas posiciones; el evento neutral
forma parte de la cronología y de su registro. Las direcciones cardinales
determinan los vectores \(\nu_d\), mientras la firma de fase determina
cuál de esos vectores se aplica.

El levantamiento entero conserva el cruce de frontera. Si
\(s:X_9\to D_9^2\) elige los representantes positivos y
\(x'=x+a\nu_d\), con \(a\in\{0,1\}\), se define
\begin{equation}
 b_d(x,a)=\frac{s(x)+a\nu_d-s(x')}{9}\in\ZZ^2.
 \label{tpk-int:borde}
\end{equation}
La divisibilidad por \(9\) se sigue de la igualdad en \(X_9\).
Si \(z\in\ZZ^2\) registra los cruces anteriores, la actualización es
\(z'=z+b_d(x,a)\). De este modo,
\[
 s(x')+9z'=s(x)+9z+a\nu_d.
\]
El dato entero de transporte permanece reconstruible después de reducir
la posición sobre el toro.

\begin{proposition}[Conservación del desplazamiento levantado]
\label{tpk-int:prop-desplazamiento}
Para una trayectoria cardinal de \(n\) pasos, cuyos indicadores de
activación son \(a_t\), se cumple
\[
 s(x_n)+9z_n=s(x_0)+9z_0+
 \sum_{t=0}^{n-1}a_t\nu_{d_t}.
\]
La misma igualdad se conserva al concatenar dos trayectorias compatibles.
\end{proposition}
\begin{proof}
La identidad de un paso es \eqref{tpk-int:borde}. Al sumarla para los
\(n\) pasos, las coordenadas intermedias se cancelan. Dividir el intervalo
de suma en dos tramos produce exactamente la ley de concatenación. En una
trayectoria cerrada, la diferencia \(z_n-z_0\) es el grado de
enrollamiento definido previamente.
\end{proof}

El mismo principio rige las evaluaciones de APP. Si una hoja toma el
valor entero \(N_t\), su registro conserva simultáneamente
\(N_t=R_t+9q_t\). La diferencia entre dos pasos verifica
\[
 N_{t+1}-N_t=(R_{t+1}-R_t)+9(q_{t+1}-q_t).
\]
La memoria de cruce espacial \(z_t\) y el cociente aritmético \(q_t\)
tienen definiciones diferentes; su conservación conjunta mantiene la
procedencia geométrica y aritmética de cada transición.

\subsubsection{Cronología de dos cursores y retornos compuestos}
\label{tpk-int:cronologia}

La involución cardinal \(\jmath\) intercambia \(N\) con \(S\) y
\(E\) con \(O\). Tras aplicar \eqref{tpk-int:cursores}, ambas
direcciones se invierten cuando finaliza un bloque de \(27\) pasos.
Escribiendo
\[
 h(\theta)=\mathbf1_{\{0,27,54,81\}}
             ((\bar\theta+1)\bmod108),
\]
la regla completa sobre la proyección observable es
\begin{equation}
 F_{\rm obs}(q)=
 \bigl(x_+',\jmath^{h(\theta)}d_+,
       x_\times',\jmath^{h(\theta)}d_\times,
       \theta+1\bigr).
 \label{tpk-int:fobs}
\end{equation}
La lectura que forma la emisión utiliza las posiciones anteriores al
transporte, según la recurrencia explícita de la
sección~\ref{sec:extension}. La actualización observable y la emisión
se componen, por tanto, con un orden fijado.

\begin{proposition}[Reversibilidad y periodo del calendario observable]
\label{tpk-int:periodo}
La aplicación \(F_{\rm obs}\) es una permutación de
\(\mathcal Q_{\rm obs}\) y satisface \(F_{\rm obs}^{108}=I\).
En un bloque de \(54\) pasos se restauran posiciones y direcciones;
la coordenada de \(\Theta_{108}\) avanza \(54\).
\end{proposition}
\begin{proof}
Dado el estado final, se recupera primero la fase anterior restando uno
en \(\Theta_{108}\). Esa fase determina si se invirtieron las
direcciones y qué cursor se desplazó. Aplicar la involución correspondiente
y restar el vector cardinal recupera el estado inicial de manera única.

Cada intervalo de \(27\) pasos contiene nueve activaciones de cada
cursor. El bloque siguiente conserva las activaciones e invierte las
direcciones, de modo que los desplazamientos enteros se cancelan en
\(54\) pasos. La regla tiene periodo \(54\) en sus incrementos de
posición y orientación; la misma cancelación vale con cualquier origen
del calendario. Tras \(108\) pasos se restaura también \(\theta\).
\end{proof}

Se reserva
\[
 \mathcal K_{27}=F_{\rm obs}^{27}:
 \mathcal Q_{\rm obs}\longrightarrow\mathcal Q_{\rm obs}
\]
para el iterado de \(27\) transiciones. Su cuarto iterado es la
identidad sobre el estado observable definido en \eqref{tpk-int:qobs}.
La proyección de la historia al calendario permite efectuar esta
comprobación finita. El registro de la trayectoria, sus evaluaciones y
los nuevos prefijos se conservan en el estado dinámico sobre el que
actúa \(\mathcal U_t\); la igualdad observable no los reinicializa.

\subsubsection{Transporte afín de la memoria y composición de estados}
\label{tpk-int:memoria}

La proyección observable admite fibras de datos aritméticos y geométricos.
Sobre una configuración \(q\), escribimos un estado en la forma
\[
 x=(q,o,h,c,\sigma,C,b,\lambda).
\]
Aquí \(o\) conserva la orientación y la hoja; \(h\), los residuos
y cocientes; \(c\), la memoria aditiva de transporte; \(\sigma\),
la firma de frontera; \(C\), el cilindro de precisión; \(b\), su
etiqueta de frontera; y \(\lambda\), la trayectoria registrada. La
firma está determinada por una aplicación
\[
 \operatorname{Sig}_q:
 \mathsf H_q\times\mathsf C_q\times\mathsf{Hist}_q
 \longrightarrow\mathsf S_q,
 \qquad \sigma=\operatorname{Sig}_q(h,c,\lambda).
\]
Los símbolos \(\mathsf H_q,\mathsf C_q,\mathsf{Hist}_q\) y
\(\mathsf S_q\) denotan, respectivamente, los conjuntos de datos
residuo--cociente, el grupo abeliano de memoria de transporte, las
trayectorias registradas y las firmas de frontera.
Las realizaciones de esta firma se especifican con sus componentes en
las construcciones regional y terminal. Esta dependencia impide tratar
la firma como un parámetro libre una vez fijados los datos que la producen.

La carta elemental de residuo y cociente es explícita. Para
\(n\in\ZZ\), sean
\[
 r=1+((n-1)\bmod9),\qquad u=\frac{n-r}{9}.
\]
Entonces \(n=r+9u\), con \(r\in D_9\) y \(u\in\ZZ\).
La unicidad se sigue de que dos representantes de \(D_9\) congruentes
módulo \(9\) son iguales. Esta carta proporciona la reconstrucción
entera local. Las torres de precisión incorporan, además, las aplicaciones
de truncamiento y los cilindros desarrollados en
\ref{sec:extension} y \ref{sec:generacion}.

Sea \(\gamma:q\to q'\) una trayectoria observable. El transporte
de las fibras se expresa mediante aplicaciones
\[
 \mathsf H(\gamma):\mathsf H_q\to\mathsf H_{q'},\qquad
 \mathsf C(\gamma):\mathsf C_q\to\mathsf C_{q'},\qquad
 \mathsf{Hist}(\gamma):\mathsf{Hist}_q\to\mathsf{Hist}_{q'}.
\]
Las aplicaciones respetan identidades y composición, y
\(\mathsf C(\gamma)\) es un homomorfismo de grupos abelianos.
El incremento \(k_{\rm tr}(\gamma)\in\mathsf C_{q'}\) cumple
\begin{equation}
 \begin{split}
 k_{\rm tr}(\operatorname{id}_q)&=0,\\
 k_{\rm tr}(\gamma_2\circ\gamma_1)
 &=k_{\rm tr}(\gamma_2)+
   \mathsf C(\gamma_2)k_{\rm tr}(\gamma_1).
 \end{split}
 \label{tpk-int:cociclo}
\end{equation}
El símbolo \(k_{\rm tr}\) designa aquí el cociclo de transporte;
la coordenada racional \(\kappa\) del registro \(K\) conserva
su notación en el capítulo de la clausura dodecafásica.

Las coordenadas transportables se actualizan por
\begin{equation}
 \begin{aligned}
 h'&=\mathsf H(\gamma)h,&
 c'&=\mathsf C(\gamma)c+k_{\rm tr}(\gamma),\\
 \lambda'&=\gamma\circ\lambda,&
 \sigma'&=\operatorname{Sig}_{q'}(h',c',\lambda').
 \end{aligned}
 \label{tpk-int:accion-memoria}
\end{equation}
La segunda ecuación es una acción afín: el término lineal transporta
la memoria anterior y el cociclo añade el incremento producido por la
trayectoria. En la realización de cruce espacial de cada cursor por separado,
\(\mathsf C(\gamma)=I\) y \(k_{\rm tr}(\gamma)\) es la suma
de los vectores de \eqref{tpk-int:borde}. Esta instancia concreta
vincula la ley abstracta con las operaciones cardinales ya definidas.

\begin{proposition}[Compatibilidad de la actualización con la concatenación]
\label{tpk-int:composicion-memoria}
La actualización \eqref{tpk-int:accion-memoria} por
\(\gamma_1\), seguida de la actualización por \(\gamma_2\),
coincide con la actualización por \(\gamma_2\circ\gamma_1\).
\end{proposition}
\begin{proof}
Para la coordenada afín, dos actualizaciones dan
\[
 c''=\mathsf C(\gamma_2)\mathsf C(\gamma_1)c
       +\mathsf C(\gamma_2)k_{\rm tr}(\gamma_1)
       +k_{\rm tr}(\gamma_2).
\]
La functorialidad de \(\mathsf C\) y
\eqref{tpk-int:cociclo} convierten esta expresión en
\(\mathsf C(\gamma_2\circ\gamma_1)c+
k_{\rm tr}(\gamma_2\circ\gamma_1)\).
La afirmación para \(h\) se sigue de la composición de
\(\mathsf H\), y para \(\lambda\), de la asociatividad de las
trayectorias. La firma final coincide porque se evalúa sobre las mismas
tres coordenadas reconstruidas.
\end{proof}

El cilindro y su frontera completan esta ley mediante las relaciones
\(\operatorname{Ext}_\gamma\), definidas en
\ref{tpk-int:bidireccionalidad}. Sus elementos determinan cuáles de
las actualizaciones anteriores conservan una prolongación admisible.
Los estados, junto con las ternas \((\gamma,x,x')\) admitidas por
esas relaciones, forman una categoría sobre las trayectorias observables:
\[
 (\gamma_2,x',x'')\circ(\gamma_1,x,x')
     =(\gamma_2\circ\gamma_1,x,x'').
\]
La identidad de \(x\) es \((\operatorname{id}_q,x,x)\).
La composición conserva a la vez el transporte aritmético y la
prolongabilidad de frontera. Esta es la estructura en la que se efectúan
los retornos con memoria y los refinamientos sucesivos.

\subsubsection{Registro dodecafásico, fase y memoria acumulada}
\label{tpk-int:dodecafase}

Las \(108\) transiciones se distribuyen en las ventanas ordenadas
\[
 E_m=\{9m-8,\ldots,9m\},\qquad 1\leq m\leq12.
\]
El registro conserva el origen de fase, la orientación y los datos
producidos durante cada ventana. Su aplicación es
\[
 \mathcal R_{12}(x)=
 \bigl((E_m,\tau_m,\sigma_m,\kappa_m,\Lambda_m)\bigr)_{m=1}^{12},
\]
donde \(\tau_m\) es la subruta de la ventana, \(\sigma_m\) su
orientación, \(\kappa_m\) el incremento entero de frontera y
\(\Lambda_m\) el registro local de incidencias. Se adopta la carta
que se desarrollará en \ref{subsec:k-registro-integral}; sus índices
de ventana corresponden a \(m=\beta+1\).
El dominio comprende las historias compatibles
del TPK y el codominio, sus registros temporales orientados. La
composición de segmentos y sus incrementos se rige por
\eqref{tpk-int:accion-memoria}. El grupo \(C_{12}\) describe, en
cambio, la traslación del índice de ventana; \(\Theta_{108}\) conserva
el calendario de pasos. Esta distinción permite precisar qué se transforma
y qué se observa al completar un ciclo.

En las coordenadas \(\theta=9b+r\), con \(0\leq r<9\), el avance
de \(9\) pasos conserva \(r\) y aumenta \(b\) en una unidad.
La adición de fases introduce el cociente
\[
 c_0(r,r')=\left\lfloor\frac{r+r'}9\right\rfloor,
\]
cuya identidad de cociclo asegura la asociatividad de la composición.
En el calendario finito se conserva \([b]_{12}\); en un registro de
vueltas se conserva \(b\in\ZZ\), o \(b\in\mathbb N_0\) para
una historia de origen fijado. La sección~\ref{sec:extension} demuestra
la sucesión exacta no escindida \eqref{eq:extension108}. Su función es
describir la relación entre estas coordenadas, antes de emplearlas en la
construcción del estado firmado.

El registro firmado recibe los datos de fase y memoria mediante la
composición que se desarrolla en la sección~\ref{sec:k-registro}; sobre la
subred compatible, las notaciones \(\mathcal S_{12}\) y
\(R_{\mathrm{sgn}}\) designan el mismo lector:
\begin{equation}
 \mathcal S_{12}\equiv R_{\mathrm{sgn}}
 =\Pi_H\circ\mathcal E_{108}^{90,120}
                         \circ\mathcal R_{12}.
 \label{tpk-int:registro-causal}
\end{equation}
Su dominio es el estado terminal con memoria; su salida pertenece a
\(\ZZ^{12}\). La transformación reversible entre ese estado firmado
y \(K\), así como la lectura racional de \(K\), se demuestran en
el capítulo correspondiente. Las operaciones de registro anteceden a
esas transformaciones. La reconstrucción de un registro desde sus
diferencias cíclicas verifica su recuperabilidad; la procedencia dinámica
debe conservarse mediante los eventos que lo produjeron.
La ecuación~\eqref{tpk-int:registro-causal} fija el tipo y la composición
efectiva. El estado terminal es producido por las ciento ocho actualizaciones
\(\mathcal U_t=\mathsf{Upd}_t\circ\mathsf{Tra}_t\circ\mathsf{Sel}_t\);
\(\mathcal R_{12}\) conserva las doce ventanas con orientación, acarreo,
frontera y ledger, y \(\mathcal E_{108}^{90,120}\) publica los canales
\((b^{90},b^{120},Q_{\rm TPK})\). La
sección~\ref{sec:k-registro} incorpora la evaluación exacta y demuestra que
\(\Pi_H\) produce \(U_{\rm term}\) antes de la inversión de Hadamard que
produce \(K\). El verificador focal comienza deliberadamente en la coordenada
de canales para comprobar de manera focal y exacta la carta posterior y su
retorno; no se confunde el alcance de ese programa auxiliar con el alcance de
la construcción matemática.

Para la traza codificada \((d_t)\) conservada por el registro, el
lector de eventos cuenta los códigos \(\{2,4,6,8\}\) de cada
ventana con la orientación
\((+,+,+,-,-,-,+,+,+,-,-,-)\). Estos códigos y los símbolos
cardinales de \(\mathcal D\) son coordenadas distintas del
registro. El capítulo~\ref{sec:k-registro} desarrolla este lector,
la descomposición integral \(1+5+6\), las congruencias que
caracterizan su imagen y su inversa exacta. La prueba de inversión
establece qué datos conserva esa carta; la definición de
\(\mathcal R_{12}\) especifica la historia orientada sobre la que
actúa.

\subsubsection{Incidencia de fases y construcción interna de los canales}
\label{tpk-int:incidencia}

El selector determina tres subconjuntos de \(D_9\):
\[
 H_+=\{1,4,7\},\qquad H_-=\{2,5,8\},\qquad H_0=\{3,6,9\}.
\]
Sea \(H_{\rm act}=H_+\sqcup H_-\) el conjunto de fases activas y
sea \(O_{\rm ph}=D_9\setminus\{9\}\) la carta anterior al retorno
marcado. Denotemos por \(P_2(H_{\rm act})\) el conjunto de pares
no ordenados de fases activas. La incidencia interna está dada por
\begin{equation}
 \mathcal F^{\rm ph}_6=H_{\rm act}\times P_2(H_{\rm act}),\qquad
 \mathcal F^{\rm ph}_8=O_{\rm ph}\times P_2(H_{\rm act}),\qquad
 \Theta_{\rm ph}=D_9\times P_2(H_{\rm act}).
 \label{tpk-int:fases90120}
\end{equation}
Los índices describen la cantidad de posiciones de la primera componente.
La construcción utiliza las fases y el origen del calendario ya fijados;
las realizaciones excepcionales posteriores transportarán estas
incidencias a sus propios dominios.

\begin{proposition}[Cardinales e incidencia orientada de fase]
\label{tpk-int:cardinales}
Los conjuntos de \eqref{tpk-int:fases90120} tienen cardinales
\(90,120,135\), respectivamente. Si
\[
 \partial_{\rm ph}=\Theta_{\rm ph}\times\{-1,+1\},\qquad
 E_{\rm ph}=\partial_{\rm ph}\sqcup\{\infty\},\qquad
 L_{\rm ph}=H_0\times P_2(H_{\rm act}),
\]
entonces
\[
 |\partial_{\rm ph}|=270,\quad |E_{\rm ph}|=271,\quad
 |L_{\rm ph}|=45,\quad
 |\partial_{\rm ph}\sqcup L_{\rm ph}|=315.
\]
Trasladar simultáneamente el origen de fase y el retorno marcado conserva
todos estos cardinales y sus inclusiones.
\end{proposition}
\begin{proof}
Hay seis fases activas, por lo que
\(|P_2(H_{\rm act})|=\binom62=15\). Los productos por
\(6,8,9\) dan \(90,120,135\). La orientación duplica \(135\);
la marca de retorno añade un elemento; las tres fases neutrales dan
\(3\cdot15=45\). Una traslación del calendario transporta cada
factor por una biyección y lleva el punto marcado a su imagen. Los
productos y las uniones disjuntas conservan así la incidencia.
\end{proof}

El cardinal \(271\) tiene aquí una realización por incidencia de
fase. La completación cúbica de la sección~\ref{sec:extension} tiene
su propia construcción de la corona \(1000-729\). La correspondencia
entre ambas realizaciones debe conservar los mapas y las marcas que las
identifican, además del cardinal. Análogamente, el selector
\(\varepsilon_9\), la incidencia \((90,120)\) y el extractor
\(\mathcal E_{108}^{90,120}\) conservan los dominios y operaciones
que se acaban de distinguir.

\subsubsection{Inversión de trayectorias y prolongaciones conjugadas}
\label{tpk-int:bidireccionalidad}

La bidireccionalidad comienza en una operación sobre trayectorias, antes
de asignarles valores escalares. Para
\(\gamma=(x_0;d_1,\ldots,d_n)\), con extremo \(x_n\), definimos
\[
 \operatorname{rev}\gamma
   =(x_n;\jmath d_n,\ldots,\jmath d_1).
\]
El origen de la trayectoria inversa es el extremo de la original.
La composición satisface
\[
 \operatorname{rev}^2\gamma=\gamma,\qquad
 \operatorname{rev}(\gamma_2\circ\gamma_1)
   =\operatorname{rev}\gamma_1\circ\operatorname{rev}\gamma_2,
\]
y el desplazamiento entero cambia de signo. Estas identidades se obtienen
invirtiendo el orden de los pasos y sustituyendo cada vector cardinal por
su opuesto. En un lazo se sigue
\(\operatorname{wind}(\operatorname{rev}\gamma)
=-\operatorname{wind}(\gamma)\).

Sobre los registros, la inversión requiere transportar asimismo el orden
de sus componentes y la orientación de sus eventos. La inversión de
ruta y una conjugación de firma son operaciones distinguibles: su
composición se especifica cuando se aplica a un registro concreto. En
particular, la involución regional que intercambia las componentes de
clausura y propagación conserva el eje de autoescala; su fórmula se
presenta en la sección~\ref{mono-ampl:regiones}. Esta operación regional
actúa sobre otro dominio que la inversión de una trayectoria cardinal.

Las extensiones sobre una ruta \(r:q\to q'\) se expresan mediante
relaciones entre fibras de estados compatibles. Si \(\mathcal F_q\)
denota la fibra de estados completos sobre \(q\), se exige
\[
 \operatorname{Ext}_r\subseteq\mathcal F_q\times\mathcal F_{q'},
 \qquad \Delta_{\mathcal F_q}\subseteq
        \operatorname{Ext}_{\operatorname{id}_q},
\]
\[
 \operatorname{Ext}_{r_2}\circ\operatorname{Ext}_{r_1}
       \subseteq\operatorname{Ext}_{r_2\circ r_1}.
\]
Su composición conserva las evaluaciones de ambas hojas, la orientación,
los cocientes, el prefijo y la frontera. La inversión de la ruta no
implica por sí sola que toda relación de extensión sea una biyección.
La prolongación conjugada se define en el dominio donde el transporte de
esos datos satisface las compatibilidades correspondientes.

\subsubsection{Holonomía nonádica y continuidad del refinamiento}
\label{tpk-int:holonomia}

Una vuelta de fase compone las transiciones sobre el estado con memoria:
\begin{equation}
 \Gamma_{9,k}=\mathcal U_{k+8}\circ\cdots\circ\mathcal U_k.
 \label{tpk-int:gamma}
\end{equation}
Sobre una historia individualizada, esta composición funcional realiza
la relación de holonomía; en el dominio completo se conserva
\(\Gamma_{9,k}\subseteq\widehat X_k\times\widehat X_{k+9}\),
con las prolongaciones admisibles desarrolladas más adelante.
La base cíclica tiene nueve fases, mientras
la fibra conserva la trayectoria, los cocientes y la frontera. La
monodromía es la acción de una vuelta de esa base sobre las fibras; la
holonomía \(\Gamma_{9,k}\) expresa su realización por las transiciones
efectivas del TPK.

Si \(g\) observa la fase y \(M\) el número de vueltas conservadas,
\[
 g(\Gamma_{9,k}x)=g(x),\qquad M(\Gamma_{9,k}x)=M(x)+1.
\]
La primera identidad expresa clausura de fase y la segunda identifica
la nueva profundidad temporal. La prolongación coinductiva conserva
además sus cilindros y prefijos. Las restricciones entre profundidades
componen un sistema inverso; una sección compatible reúne sus
antecedentes sin reemplazarlos por la última observación. El desarrollo
de esta aritmética de historias y de su unidad ocupa
\ref{hist:seccion}--\ref{hist:doble-varianza}; la construcción de las
regiones y sus lectores se realiza en la sección~\ref{sec:generacion}.

La transferencia de fase sobre emisiones ya construidas tiene otro
dominio. Sean \(\mathcal A_+\) y \(\mathcal A_\times\) los
alfabetos de firmas aditivas y multiplicativas, de cardinales \(18\)
y \(26\). Su producto identifica el catálogo \(\mathcal U_6\).
Para \(f:\mathcal A_+\times\mathcal A_\times\to\RR\),
las medias de hoja son
\[
 (E_+f)(s,e)=\frac1{18}\sum_{s'\in\mathcal A_+}f(s',e),\qquad
 (E_\times f)(s,e)=\frac1{26}\sum_{e'\in\mathcal A_\times}f(s,e').
\]
La construcción de estas emisiones y sus cardinales se efectúa mediante
la recurrencia de la sección~\ref{sec:extension}. Una vez obtenidas,
el operador de transferencia se define sobre
\(\mathcal U_6\times C_9\) por
\[
 P_+=E_+,\quad P_-=E_\times,\quad P_0=I,\qquad
 \mathcal K_{\rm ph}((u,r),(v,r+1))=P_{M_{\rm ph}(r)}(u,v),
\]
con peso cero para las restantes transiciones de fase. El índice
\(r\) utiliza el representante positivo de la clase cíclica.
El orden de fases \((+1,-1,0)^3\) produce entonces la renovación
\[
 \mathcal K_{\rm ph}^{9}\big|_r=E_+E_\times.
\]
En efecto, ambas medias son idempotentes, conmutan y promedian factores
diferentes; su composición asigna peso \(1/(18\cdot26)=1/468\)
a cada emisión. Esta igualdad caracteriza un observable de transferencia
posterior. El transporte de las historias sigue dado por
\eqref{tpk-int:gamma}, con su memoria y su frontera.

La jerarquía resultante determina el orden causal de las construcciones
posteriores.
La emisión finita construye el dominio y las regiones; las relaciones de
extensión conservan sus prefijos; la conexión nonádica los prolonga; el
registro dodecafásico determina los datos que intervienen en \(K\) y
en la clausura de \(\alpha\). La realización numérica y la realización
incidencial reciben esa estructura anterior. Este orden permite exponer
cada consecuencia en su capítulo conservando, desde el núcleo formal,
los objetos que la hacen posible.
